\documentclass[conference]{../../../../Conference-LaTeX-template_10-17-19/IEEEtran}
\usepackage{cite}
\usepackage{amsmath,amssymb,amsfonts}
\usepackage{array}
\usepackage{booktabs}
\usepackage{tabularx}
\usepackage{url}
\newcolumntype{Y}{>{\raggedright\arraybackslash}X}
\renewcommand{\arraystretch}{1.12}
\title{Can Macro- and Microworlds Be Simulated Accurately? A Scale-Resolved Credibility Research Proposal}
\author{\IEEEauthorblockN{Anonymous Research Proposal}}
\begin{document}
\maketitle

\begin{abstract}
Scientific simulation spans continuum flows, climate and materials models, atomistic dynamics, electronic structure, and quantum many-body systems. The common question ``can we accurately simulate the macro- and microworld?'' has no unconditional yes-or-no answer. Accuracy is a claim about a declared quantity of interest, physical model, parameter regime, numerical or algorithmic approximation, reference evidence, uncertainty interval, and resource budget. This proposal develops the \emph{Scale-Resolved Simulation Credibility Ledger} (SRCL), a source-bounded protocol for evaluating such claims across macro continuum, atomistic, classically simulated quantum, noisy intermediate-scale quantum (NISQ), and fault-tolerant projection pathways. SRCL separates verification from validation, records a structured error budget, requires a compatible classical baseline for every quantum candidate, and prevents future hardware assumptions from being stated as observed performance. The central hypothesis is that a shared quantity-of-interest, tolerance, evidence, and cost ledger will reveal whether a scenario is numerically verified, physically validated for a stated regime, currently classically preferable, a quantum candidate for specialist benchmark study, or simply inconclusive. The protocol is design-only: it runs no HPC or quantum-hardware jobs, gathers no new data, and claims no quantum advantage. Its contribution is a falsifiable decision language that resists two misconceptions at once: macroscopic simulation is not automatically easy, and quantum computing is not automatically necessary or accurate for microscopic simulation.
\end{abstract}

\begin{IEEEkeywords}
scientific computing, verification and validation, uncertainty quantification, multiscale modeling, quantum simulation, quantum computing, error mitigation, fault tolerance
\end{IEEEkeywords}

\section{Introduction}
Simulation is among science’s most powerful ways of connecting a theory to a quantity that cannot be directly observed in every condition. A weather or fluid model can forecast a field over a domain. A molecular model can estimate a reaction-energy difference. A quantum many-body model can predict a correlation function. A quantum computer can, in principle, represent the state evolution of a quantum system using quantum degrees of freedom. These achievements can create the impression that the macroworld is straightforward to simulate while only the microworld demands powerful quantum hardware. The impression is incomplete.

Macroscopic systems can be hard because a governing law must be closed, discretized, calibrated, coupled to unresolved scales, and evaluated under uncertain inputs or chaotic sensitivity. A microscale problem can be tractable classically when a good electronic-structure approximation, Monte Carlo method, tensor network, reduced model, or high-performance solver matches its requested observable and precision. Conversely, a strongly correlated or real-time quantum problem can overwhelm a classical representation even when the physical sample is microscopic. Scale alone does not choose the algorithm or determine credibility.

The correct question is therefore conditional: \emph{under what model, numerical, data, algorithmic, and hardware-error conditions is a simulation credible for its declared quantity of interest?} Verification-and-validation methodology separates whether a code and numerical solution solve the mathematical model from whether that model represents a physical target in the stated regime \cite{oberkampf_trucano2002,oberkampf_trucano2007}. This distinction remains essential when the computational pathway changes from a classical PDE solver to a quantum circuit. A fast computation, a small residual, a high qubit count, or an attractive benchmark alone cannot validate the underlying physics.

Quantum information science provides a legitimate motivation for reexamining difficult microscopic simulation. NIST describes the field as joining quantum physics with information theory and notes that quantum computers may, in theory, simulate quantum matter and tackle certain presently unsolved problems \cite{nist_qis_page}. Foundational proposals explain why a controllable quantum system can simulate another quantum system \cite{feynman1982,lloyd1996}. Yet the application-level question remains: after Hamiltonian mapping, state preparation, algorithmic approximation, measurement sampling, noise, error mitigation or correction, and resource accounting, does the candidate meet the same quantity-of-interest tolerance as an appropriate classical comparator?

This paper presents the \emph{Scale-Resolved Simulation Credibility Ledger} (SRCL). SRCL is not a new physical law, a benchmark result, or a claim that an existing quantum platform has surpassed a named classical method. It is a research design that makes the required evidence explicit. A scenario enters the ledger with a quantity of interest, units, tolerance, model regime, error terms, reference/validation route, and cost interval. It receives a conditional status only after those fields and a compatible baseline have been reviewed. The intended outcome is a more precise scientific statement than either ``we simulated it'' or ``a quantum computer will solve it.''

\section{What Accuracy Means in a Scientific Simulation}
\subsection{Quantity of interest, model, and reference}
Let $Q^\star$ denote the physical quantity of interest and $Q_h$ a computed estimate under a declared model and computational representation. The physical discrepancy is
\begin{equation}
E_Q=Q_h-Q^\star.
\label{eq:physical-error}
\end{equation}
In most scientific applications $Q^\star$ is not directly known. A credible claim therefore cannot report only a numerical difference; it must state the reference route: an analytic solution, a manufactured solution, a numerically refined reference, a controlled experiment, an observational dataset with uncertainty, or an explicitly unvalidated extrapolation. The reference has its own uncertainty and applicability boundary.

SRCL records a conservative structured error budget rather than assuming a universal independent-error formula:
\begin{equation}
\begin{aligned}
\epsilon_Q^{\mathrm{claim}}\ \lesssim\ &\epsilon_{\mathrm{model}}+\epsilon_{\mathrm{input}}+\epsilon_{\mathrm{repr}}+\epsilon_{\mathrm{algorithm}}\\
&+\epsilon_{\mathrm{sampling}}+\epsilon_{\mathrm{coupling}}+\epsilon_{\mathrm{data}}+\epsilon_{\mathrm{quantum}}.
\end{aligned}
\label{eq:error-ledger}
\end{equation}
The symbol $\lesssim$ in \eqref{eq:error-ledger} is intentional. Terms can be correlated, bounded asymmetrically, or assessed through an ensemble rather than summed. The ledger is an auditable decomposition, not an assertion that errors add exactly. For a purely classical card, $\epsilon_{\mathrm{quantum}}=0$. For a quantum card, that term expands into mapping, state-preparation, circuit/algorithm, compilation, physical-noise, mitigation or correction, and measurement components.

The first research discipline is to define a tolerance $\tau_Q$ before selecting a platform. A simulation may be adequate for the sign of a phase change but inadequate for a reaction barrier, a long-time trajectory, or a safety-relevant load. It can be verified with respect to a mathematical model but not validated for an unmeasured regime. It can be statistically precise but systematically biased by a model discrepancy. These categories are not rhetorical distinctions; they control which error terms must be closed.

\subsection{Verification, validation, and uncertainty}
Verification asks whether a numerical implementation solves the stated mathematical equations to the claimed accuracy. It may involve code tests, conservation checks, manufactured solutions, convergence studies, or independently generated high-accuracy references. Validation asks whether the model represents the physical target for the intended conditions. It requires an appropriately related physical reference, measurement uncertainty, calibration policy, and extrapolation assessment. Verification does not validate a closure model; validation evidence does not excuse a nonconverged numerical solution \cite{oberkampf_trucano2002,oberkampf_trucano2007}.

Uncertainty quantification complements both activities. Inputs can be uncertain, model forms incomplete, numerical approximations finite, and data noisy. Multifidelity methods are useful precisely because low-cost and high-cost models can be combined while a high-fidelity model remains connected to the accuracy or convergence argument \cite{peherstorfer2018}. SRCL adopts this logic across scales: a cheap surrogate, a coarse mesh, a reduced-order model, or a shallow quantum circuit may support exploration, but it cannot silently inherit the credibility of a separate reference calculation.

\section{Evidence Survey Across Scales and Pathways}
\subsection{Macroscopic and multiscale simulation}
Continuum and macroscopic models often benefit from conservation laws, scale separation, and mature numerical methods. Their success is real but conditional. A continuum assumption can fail near interfaces or rarefied conditions. Turbulence or constitutive closures can dominate the uncertainty even when a mesh converges. Coupled subsystems can introduce interface error. Chaotic systems can make trajectory agreement inappropriate while distributional or forecast metrics remain meaningful. Thus, ``macroworld'' should be an SRCL card with equations, closure, conservation diagnostics, mesh/time refinement, input uncertainty, and validation regime, not a pass-by-default category.

The same point applies to multiscale models. A macro quantity may depend on microscopic constitutive information; an atomistic or reduced calculation may provide that information with an approximation and sampling error; a coupling mechanism transfers assumptions between scales. SRCL therefore treats $\epsilon_{\mathrm{coupling}}$ as an explicit field. A simulation that has excellent component solvers but an unverified interface may remain inconclusive for the coupled quantity of interest.

\subsection{Microscopic simulation on classical computers}
The microscopic domain is not a single computational class. Classical molecular dynamics, Monte Carlo, electronic-structure methods, finite-temperature techniques, tensor networks, and embedding approaches each exploit structure in a particular problem. A reduced basis can be accurate for one observable and inaccurate for another. A classical algorithm can remain preferable when it meets the tolerance with lower total cost or a stronger reference route. The appropriate comparison is not ``classical versus quantum'' in general; it is a compatible model, observable, precision, and resource comparison.

Representation choices are central. In chemistry and materials work, an active-space or basis truncation may reduce resource demands but introduce a representation error. Techniques that exchange qubits or circuit depth for additional measurement burden illustrate the multi-dimensional nature of this trade-off \cite{takeshita2020}. SRCL records such a choice in $\epsilon_{\mathrm{repr}}$, $\epsilon_{\mathrm{sampling}}$, and the cost interval rather than reporting only one hardware metric.

\subsection{Quantum simulation: promise and constraints}
For a quantum system described by a target Hamiltonian $H$ and initial state $\rho_0$, a typical dynamical quantity can be written
\begin{equation}
Q(t)=\mathrm{Tr}\!\left[Oe^{-iHt}\rho_0e^{iHt}\right].
\label{eq:quantum-qoi}
\end{equation}
A quantum algorithm must map $H$ and $O$ to a finite qubit representation, prepare an appropriate state, approximate the evolution or eigenvalue procedure, execute and compile a circuit, estimate measurement expectations, and correct or mitigate device errors. Each step can affect the requested $Q(t)$. Universal quantum simulation provides an important theoretical framework \cite{lloyd1996}, but universality is not a certificate of a usable tolerance or resource budget.

NISQ research recognizes this distinction. Noise limits reliable circuit depth and makes mitigation strategies relevant \cite{preskill2018}. Cross-validated evidence on pair-correlated electron simulations demonstrates why mitigation can be scientifically useful while still requiring resource and scaling assessment \cite{obrien2023}. Cross-validated logical-qubit work demonstrates error-suppression progress, but a device-level logical error result is not by itself an application-level simulation validation \cite{google2023}. The ledger therefore records device characterization as one evidence type, not as a substitute for the physical reference and model-validation route.

\begin{table*}[!t]
\caption{Scale and pathway distinctions retained by SRCL. Constraint candidates are required ledger fields, not observed rankings.}
\label{tab:pathways}
\centering
\small
\begin{tabularx}{\textwidth}{lYYY}
\toprule
Model card & Primary use case & Required credibility fields & Scope safeguard \\
\midrule
C1: macro continuum/PDE & fields, flows, structures, coupled systems & governing equations, closure, conservation, mesh/time studies, validation regime & numerical convergence is not physical validation \\
C2: atomistic classical/ab-initio & molecular or material quantities & model/basis, finite-size, sampling, calibration/reference route & microscopic scale does not imply quantum hardware need \\
C3: classical quantum many-body & selected correlated/quantum observables & Hamiltonian, truncation/sampling, complexity indicator, reference & compare to a stated observable and precision \\
C4: NISQ quantum simulation & near-term algorithm and device candidate & mapping, state preparation, circuit, noise, mitigation, sampling, classical baseline & device metric is not model validation \\
C5: fault-tolerant projection & future resource scenario & logical-error assumption, resource model, projection flag, baseline & never report a projection as observed performance \\
C6: cross-scale coupling & coupled macro--micro quantity & interface variables, coupling error, consistency, validation & component accuracy does not guarantee coupled accuracy \\
\bottomrule
\end{tabularx}
\end{table*}

\section{Research Gaps and Idea Selection}
\subsection{Gap ledger}
The Survey Agent accepts two connected gaps. First, accuracy language often combines three noninterchangeable pieces of evidence: numerical verification, physical validation, and computing-platform performance. Second, macro continuum, atomistic, classically simulated quantum, NISQ, and projected fault-tolerant studies rarely share one explicit quantity-of-interest and error/cost contract. That absence makes an apparently impressive result hard to compare with a classical baseline or a different scale.

The Survey Agent rejects an appealing but unsupported proposition: that a powerful quantum computer is necessary for all microscale simulation. Classical microscopic methods often remain the most credible option for a declared task. Conversely, the existence of a hard classical problem does not make an arbitrary quantum circuit accurate. The path from a hard Hamiltonian to a useful result contains the full sequence in Eq. \eqref{eq:quantum-qoi} and the quantum error ledger in Eq. \eqref{eq:error-ledger}.

\subsection{Selected idea}
The Idea Agent explored a universal quantum-simulation score, a macro-only convergence checklist, a hardware-error dashboard, and SRCL. The universal score was rejected because it lacks a quantity of interest, tolerance, and compatible baseline. The macro checklist was rejected because it cannot represent algorithmic mapping or hardware terms. The hardware dashboard remains useful but cannot show whether a model predicts physics. SRCL is selected because it unifies the error, evidence, and cost questions without pretending that every card should use the same solver.

SRCL’s central hypothesis is: \emph{a shared quantity-of-interest, tolerance, error decomposition, reference/validation route, and resource ledger will reveal whether a simulation pathway is justified, classically preferable, a quantum candidate for specialist study, or inconclusive.} A quantum candidate may be attractive if it provides a plausible tolerance/resource case relative to a compatible classical baseline. That status is intentionally weaker than an observed speedup or demonstrated application advantage.

\begin{figure*}[!t]
\centering
\fbox{\begin{minipage}{0.94\textwidth}
\centering\textbf{Scale-Resolved Simulation Credibility Ledger (SRCL)}\\[3pt]
\small
\textit{Declared quantity of interest} $\rightarrow$ \textit{scale/pathway card} $\rightarrow$ \textit{structured error ledger} $\rightarrow$ \textit{reference and validation route} $\rightarrow$ \textit{conditional credibility status}.\\[5pt]
\begin{tabular}{c@{\quad}c@{\quad}c@{\quad}c}
macro PDE & atomistic/classical quantum & NISQ candidate & fault-tolerant projection\\
closure, mesh, coupling & basis, sampling, reference & mapping, circuit, noise & logical-error/resource assumptions
\end{tabular}\\[5pt]
\textit{Outputs:} verified, validated, classically preferable, quantum candidate advantage, inconclusive, or model invalid.
\end{minipage}}
\caption{SRCL decision flow. The editable source is retained as \texttt{figures/scale\_resolved\_credibility\_ledger.drawio}; the schematic describes a proposed workflow, not a completed benchmark.}
\label{fig:srcl}
\end{figure*}

\section{ExperimentDesign: SRCL Computational Protocol}
\subsection{Design-only study objective}
The ExperimentDesign Agent selects the computational-digital template. Its unit of analysis is one quantity-of-interest $\times$ model card $\times$ computational pathway $\times$ parameter-scenario tuple. The protocol does not run a PDE solver, submit a job to an HPC queue, invoke a cloud quantum processor, or collect validation data. Instead, it specifies the information and checks that a future benchmark must provide before an accuracy or advantage statement can be written.

The baseline card \texttt{C0} defines the reference route. \texttt{C1} records a macro continuum/PDE scenario. \texttt{C2} records an atomistic classical or ab-initio scenario. \texttt{C3} records a classically simulated quantum-many-body scenario. \texttt{C4} records a NISQ quantum candidate. \texttt{C5} records a fault-tolerant resource projection. \texttt{C6} records cross-scale coupling. Each card preserves units, model scope, tolerance, and provenance. A comparison is prohibited when the cards use different observables, different physical Hamiltonians/equations, incompatible precision, or a projected quantum cost against an observed classical cost without a clear qualifier.

\subsection{Error and evidence contract}
Table \ref{tab:contract} defines the operational fields. The procedure first verifies basic consistency: units, conservation or invariants where applicable, state/Hamiltonian validity for quantum cards, and the mathematical identity of the quantity of interest. It next asks which terms in Eq. \eqref{eq:error-ledger} have a bound, an estimate, a sensitivity range, or only a placeholder. A missing term cannot be silently set to zero. The output becomes \texttt{INCONCLUSIVE} if the evidence is incomplete, even when one component has an excellent benchmark.

\begin{table*}[!t]
\caption{SRCL data contract. All values are planned inputs, source-bounded estimates, or derived analysis fields; this proposal reports no executed simulation.}
\label{tab:contract}
\centering
\small
\begin{tabularx}{\textwidth}{lYYY}
\toprule
Category & Required fields & Validation rule & Failure response \\
\midrule
Quantity and scope & $Q$, units, tolerance $\tau_Q$, regime, prediction horizon & definition unchanged inside comparison set & \texttt{MODEL\_INVALID} \\
Model and representation & equations/Hamiltonian, closure/basis/mapping, parameters & approximation source and domain declared & \texttt{INCONCLUSIVE} \\
Numerical/algorithmic & mesh/basis, solver/algorithm, stopping or approximation bounds & refinement or bounded-error route declared & \texttt{INCONCLUSIVE} \\
Reference and validation & analytic/refined/data/device reference, uncertainty, calibration policy & verification and validation statuses separated & no physical-validation label \\
Quantum-specific & mapping, preparation, circuit, noise, mitigation/correction, samples & compatible classical baseline and projection flag present & \texttt{INCONCLUSIVE} \\
Sensitivity and provenance & ranges, source ID, code/device version, observed/proposed flag & status stability or switch boundary reported & withhold unique ranking \\
\bottomrule
\end{tabularx}
\end{table*}

\subsection{Conditional status rules}
\texttt{VERIFIED\_WITHIN\_TOLERANCE} applies when numerical or algorithmic error is bounded for the stated mathematical model and quantity. It deliberately does not mean that nature has been validated. \texttt{VALIDATED\_FOR\_DECLARED\_REGIME} applies when a suitable physical comparison supports the model in its stated domain with an uncertainty range. \texttt{CLASSICALLY\_PREFERABLE} applies when a compatible classical card meets the tolerance at lower declared cost or risk. \texttt{QUANTUM\_CANDIDATE\_ADVANTAGE} applies only when a quantum card has a complete mapping, error, sampling, and resource ledger relative to a classical baseline; it authorizes a specialist-reviewed benchmark, not a public advantage claim. \texttt{INCONCLUSIVE} and \texttt{MODEL\_INVALID} protect against ranking incomplete or incompatible scenarios.

Sensitivity analysis varies mesh/basis resolution, model parameters, solver tolerance, sample count, mapping choice, noise proxy, mitigation method, and logical-error assumption over declared intervals. A status that changes across the interval must report the switch variable and cannot be advertised as a general conclusion. This rule is particularly important for fault-tolerant projections, where logical-qubit and resource assumptions are necessarily prospective.

\section{Expected Outcome Branches and Cross-Scale Bridge}
The Author stage is limited to the conditional language in Table \ref{tab:branches}. SRCL does not predict a winner. Its most valuable result may be an \texttt{INCONCLUSIVE} label that identifies exactly which reference, model, or resource term is missing. Such a result is scientifically useful because it converts an ill-posed “can we simulate it?” question into a targeted evidence request.

\begin{table}[!t]
\caption{Conditional result branches. These are planned interpretations, not observed simulation outcomes.}
\label{tab:branches}
\centering
\footnotesize
\begin{tabularx}{\columnwidth}{lY}
\toprule
Label & Authorized conditional interpretation \\
\midrule
\texttt{VERIFIED\_WITHIN\_TOLERANCE} & The implementation pathway is numerically or algorithmically controlled for the stated mathematical model and $Q$. \\
\texttt{VALIDATED\_FOR\_DECLARED\_REGIME} & A physical comparison supports conditional credibility for the stated regime; no extrapolation claim is implied. \\
\texttt{CLASSICALLY\_PREFERABLE} & The compatible classical pathway currently closes the ledger at lower declared cost or risk. \\
\texttt{QUANTUM\_CANDIDATE\_ADVANTAGE} & A complete quantum ledger justifies a specialist-reviewed benchmark; no observed advantage is claimed. \\
\texttt{INCONCLUSIVE} & Missing baseline, error, validation, or sensitivity evidence prevents an accuracy or advantage ranking. \\
\texttt{MODEL\_INVALID} & Units, scope, conservation, provenance, or observed/projected-status checks fail. \\
\bottomrule
\end{tabularx}
\end{table}

The cross-scale bridge is the quantity of interest, not an arbitrary computational label. A macroscopic force, a molecular free-energy difference, an electronic correlation, and a quantum correlation function require different model cards, but each can state units, tolerance, reference route, and error terms. This permits an honest comparison of evidence without flattening the physics. A quantum simulation becomes valuable when it changes the credible attainable set for a specific $Q$, not merely when it represents a microscopic system.

The bridge also makes multifidelity design more disciplined. A low-fidelity macro surrogate can screen parameters while a high-fidelity calculation establishes convergence. A classical approximate quantum model can generate an initial state or comparison reference for a quantum circuit. A quantum subproblem can be embedded into a larger classical workflow. These architectures can be productive, but SRCL requires interface variables and coupling error to be stated. It rejects an implicit assumption that a locally accurate component automatically makes a global coupled prediction accurate.

\section{Research Implications and Validation Hierarchy}
\subsection{A hierarchy for simulation claims}
SRCL distinguishes five escalating statements. A computation first produces an output. It becomes a numerical result only after code/solution verification. It becomes a model prediction only after the model scope is declared. It becomes a physically credible prediction only after validation evidence supports the relevant regime. It becomes a platform-comparison claim only after compatible accuracy and total resource accounting. This hierarchy prevents a successful circuit execution or a converged mesh from being reported as a general scientific prediction.

The hierarchy is valuable for quantum computing because it keeps promise and evidence aligned. The ability of quantum hardware to process quantum information, the theoretical universality of quantum simulation, and a demonstration of error mitigation are meaningful advances. They do not imply that every chemical, material, or many-body quantity is more accurately or cheaply obtained on an available device. Conversely, a classical method’s current success on a benchmark does not rule out a future quantum pathway for another observable, tolerance, or system-size regime. The comparison must be rebuilt for the declared tuple.

\begin{table*}[!t]
\caption{SRCL claim hierarchy. Each level inherits the requirements above it; no lower-level success licenses a higher-level conclusion.}
\label{tab:hierarchy}
\centering
\small
\begin{tabularx}{\textwidth}{lYYY}
\toprule
Claim level & Minimum evidence & Allowed statement & Prohibited escalation \\
\midrule
Computation & executable workflow and output record & ``A calculation produced this output.'' & ``The quantity is accurate.'' \\
Verification & tests, refinement or bounded-error evidence & ``The stated model/algorithm is solved within tolerance.'' & ``The physical model is validated.'' \\
Validation & related data/reference, uncertainty, regime comparison & ``The model is conditionally credible in this regime.'' & ``The model is universal or extrapolation-free.'' \\
Resource comparison & compatible $Q$, tolerance, model, cost and risk ledger & ``This pathway is currently preferable for this tuple.'' & ``The platform is universally superior.'' \\
Quantum candidate & complete mapping/noise/mitigation/logical-resource and baseline record & ``A specialist-reviewed quantum benchmark is justified.'' & ``A quantum advantage has been observed.'' \\
\bottomrule
\end{tabularx}
\end{table*}

\subsection{Practical research questions enabled}
The proposal enables targeted follow-on questions. Which macro closure term dominates a forecast quantity after mesh convergence? Which basis or active-space approximation is responsible for a microscopic observable’s error? For a quantum candidate, does state preparation, circuit approximation, sampling, physical noise, or measurement mitigation dominate the uncertainty? How does the answer change if the tolerance is relaxed or tightened? At what size or coupling strength does a classical reference cease to be practical, and what quantum resource assumptions then become decisive? These questions are sharper than a single scale-based claim and can be tested with future authorized simulations.

SRCL also supports reproducibility. Every scenario carries a ``credibility passport'': quantity definition, units, tolerance, equations/Hamiltonian, parameter domain, approximation list, error ledger, reference/validation state, cost model, source IDs, code or device version, and observed/projected flag. A later researcher can inspect whether two results are comparable, add a missing error estimate, or identify the assumption responsible for a status switch. The passport is especially important when a proposed fault-tolerant resource estimate and an observed noisy-device benchmark appear in the same discussion.

\section{Risks, Limitations, and Human Review}
SRCL cannot eliminate model discrepancy by naming it. A model may be valid only in a regime that is poorly sampled, and a reference dataset may not cover the desired extrapolation. Error decompositions can contain dependent terms and unknown correlations. The ledger therefore records an evidence type and uncertainty relation for each term rather than forcing every source into a single scalar. A small reported error bar is not credible if it omits a dominant model-form or coupling uncertainty.

Quantum-specific risks are equally concrete. A Hamiltonian mapping may discard degrees of freedom; an initial-state assumption may be weak; a circuit approximation may not meet the desired time or energy accuracy; sampling can dominate total cost; mitigation can trade bias against variance; and a fault-tolerant projection depends on logical-error and architecture assumptions. Published hardware progress and mitigation experiments are sources of bounded context, not permission to state a generic application advantage. A classical baseline can also evolve, so resource comparisons require dates, software versions, and compatible cost definitions.

Required review is specific. A domain scientist must assess physical-model validity and regime relevance. A numerical analyst must assess discretization, solver, and uncertainty evidence. A quantum-computing specialist must inspect mapping, algorithm, device, mitigation, and fault-tolerant assumptions. A bibliographic reviewer must verify DOI and publisher pages before external publication. The execution policy is \texttt{DESIGN\_ONLY}; no production or quantum-hardware action is authorized.

\section{Conclusion}
We can simulate many macro- and microscale phenomena accurately enough for declared scientific goals, but no scale receives accuracy by default. A credible simulation claim identifies its quantity of interest, model domain, tolerance, numerical/algorithmic approximation, reference route, uncertainty, and resource budget. Macroscopic systems can be limited by closure, coupling, uncertainty, and validation. Microscopic systems can often be addressed by classical methods. Quantum computing is a serious candidate for selected hard quantum problems, yet its mapping, algorithmic, sampling, noise, mitigation, correction, and resource terms must close the same credibility ledger.

The proposed SRCL framework makes this answer operational. It separates verification, validation, and platform performance; demands a compatible classical baseline for quantum candidates; and returns an inconclusive status when evidence is incomplete. The contribution is deliberately not a claim of universal simulation accuracy or observed quantum advantage. It is a falsifiable, scale-resolved method for deciding what a simulation result can honestly support and what evidence a future study must add.

\appendices
\section{Evidence Coverage and Claim Boundary}
The source registry contains twelve entries. NIST’s public quantum-information page was browser verified \cite{nist_qis_page}. Dual-engine OpenAlex/AnySearch searches cross-validated the selected verification/validation, multifidelity, NISQ, error-mitigation, logical-error, representation-accuracy, and hardware-review sources \cite{oberkampf_trucano2002,oberkampf_trucano2007,peherstorfer2018,preskill2018,obrien2023,google2023,takeshita2020,wendin2017}. Foundational quantum-simulation and chemistry-review citations require final publisher-page verification \cite{feynman1982,lloyd1996,bauer2020}. No source is used to claim that SRCL was executed, that a universal quantum advantage exists, or that a particular scientific system has been newly simulated.

\section{Symbols and Operational Definitions}
\begin{tabularx}{0.96\columnwidth}{lY}
\toprule
Symbol / term & Definition in this proposal \\
\midrule
$Q$, $Q^\star$, $Q_h$ & declared quantity of interest, physical target, and computed estimate \\
$\tau_Q$ & predeclared quantity-of-interest tolerance \\
$\epsilon_{\mathrm{model}}$ & model-form or regime discrepancy term \\
$\epsilon_{\mathrm{repr}}$ & mesh, basis, truncation, or mapping representation term \\
$\epsilon_{\mathrm{algorithm}}$ & solver, Trotter, variational, or other algorithm approximation term \\
$\epsilon_{\mathrm{quantum}}$ & mapping, preparation, circuit, noise, mitigation/correction, and measurement terms for quantum cards \\
SRCL & Scale-Resolved Simulation Credibility Ledger \\
\bottomrule
\end{tabularx}

\begin{thebibliography}{99}
\bibitem{nist_qis_page} National Institute of Standards and Technology, ``Quantum information science.'' [Online]. Available: \url{https://www.nist.gov/quantum-information-science}. Accessed: Sep. 1, 2026.
\bibitem{oberkampf_trucano2002} W. L. Oberkampf and T. G. Trucano, ``Verification and validation in computational fluid dynamics,'' U.S. DOE technical report, 2002, doi: 10.2172/793406.
\bibitem{oberkampf_trucano2007} W. L. Oberkampf and T. G. Trucano, ``Verification and validation benchmarks,'' U.S. DOE technical report, 2007, doi: 10.2172/901974.
\bibitem{peherstorfer2018} B. Peherstorfer, K. Willcox, and M. Gunzburger, ``Survey of multifidelity methods in uncertainty propagation, inference, and optimization,'' \emph{SIAM Rev.}, vol. 60, no. 3, pp. 550--591, 2018, doi: 10.1137/16M1082469.
\bibitem{feynman1982} R. P. Feynman, ``Simulating physics with computers,'' \emph{Int. J. Theor. Phys.}, vol. 21, pp. 467--488, 1982, doi: 10.1007/BF02650179.
\bibitem{lloyd1996} S. Lloyd, ``Universal quantum simulators,'' \emph{Science}, vol. 273, pp. 1073--1078, 1996, doi: 10.1126/science.273.5278.1073.
\bibitem{preskill2018} J. Preskill, ``Quantum computing in the NISQ era and beyond,'' \emph{Quantum}, vol. 2, p. 79, 2018, doi: 10.22331/q-2018-08-06-79.
\bibitem{bauer2020} B. Bauer \emph{et al.}, ``Quantum algorithms for quantum chemistry and quantum materials science,'' \emph{Chem. Rev.}, vol. 120, pp. 12685--12717, 2020, doi: 10.1021/acs.chemrev.9b00829.
\bibitem{obrien2023} T. E. O’Brien \emph{et al.}, ``Purification-based quantum error mitigation of pair-correlated electron simulations,'' \emph{Nat. Phys.}, vol. 19, pp. 1787--1792, 2023, doi: 10.1038/s41567-023-02240-y.
\bibitem{google2023} Google Quantum AI \emph{et al.}, ``Suppressing quantum errors by scaling a surface code logical qubit,'' \emph{Nature}, vol. 614, pp. 676--681, 2023, doi: 10.1038/s41586-022-05434-1.
\bibitem{takeshita2020} T. Y. Takeshita \emph{et al.}, ``Increasing the representation accuracy of quantum simulations of chemistry without extra quantum resources,'' \emph{Phys. Rev. X}, vol. 10, 011004, 2020, doi: 10.1103/PhysRevX.10.011004.
\bibitem{wendin2017} G. Wendin, ``Quantum information processing with superconducting circuits: A review,'' \emph{Rep. Prog. Phys.}, vol. 80, 106001, 2017, doi: 10.1088/1361-6633/aa7e1a.
\end{thebibliography}
\end{document}
